\documentclass[11pt]{ctexart}
\usepackage[a4paper,margin=1in]{geometry}
\usepackage{longtable,booktabs,array}
\usepackage{xcolor}
\usepackage{hyperref}
\title{Geant4-Agent 回归测试报告（中文）}
\author{自动化评测流程}
\date{2026-03-05\_231154}
\begin{document}
\maketitle
\section{统一 Pass/Fail 总表}
\begin{longtable}{p{0.06\linewidth}p{0.1\linewidth}p{0.08\linewidth}p{0.09\linewidth}p{0.08\linewidth}p{0.08\linewidth}p{0.06\linewidth}p{0.07\linewidth}p{0.28\linewidth}}
\toprule
ID & 领域 & 风格 & 参数完整性 & 期望闭环 & 实际闭环 & 缺参数 & 结果 & 失败原因 \\
\midrule
PF01 & medical & precise & full & 通过 & 通过 & 0 & 通过 & - \\
PF02 & aerospace & precise & full & 通过 & 通过 & 0 & 通过 & - \\
PF03 & industrial\_ndt & precise & full & 通过 & 通过 & 0 & 通过 & - \\
PF04 & high\_energy\_physics & precise & full & 通过 & 失败 & 3 & 失败 & expected\_complete\_but\_incomplete, unexpected\_missing\_fields \\
PF05 & nuclear\_engineering & precise & full & 通过 & 通过 & 0 & 通过 & - \\
PM01 & medical & precise & missing & 失败 & 失败 & 8 & 通过 & - \\
PM02 & aerospace & precise & missing & 失败 & 失败 & 5 & 通过 & - \\
PM03 & industrial\_ndt & precise & missing & 失败 & 失败 & 4 & 通过 & - \\
PM04 & high\_energy\_physics & precise & missing & 失败 & 失败 & 3 & 通过 & - \\
PM05 & environmental\_radiation & precise & missing & 失败 & 失败 & 10 & 通过 & - \\
FF01 & medical & fuzzy & full & 通过 & 通过 & 0 & 通过 & - \\
FF02 & aerospace & fuzzy & full & 通过 & 通过 & 0 & 通过 & - \\
FF03 & security\_screening & fuzzy & full & 通过 & 通过 & 0 & 通过 & - \\
FF04 & education & fuzzy & full & 通过 & 通过 & 0 & 通过 & - \\
FF05 & semiconductor & fuzzy & full & 通过 & 通过 & 0 & 通过 & - \\
FM01 & medical & fuzzy & missing & 失败 & 失败 & 8 & 通过 & - \\
FM02 & aerospace & fuzzy & missing & 失败 & 失败 & 12 & 通过 & - \\
FM03 & nuclear\_engineering & fuzzy & missing & 失败 & 失败 & 3 & 通过 & - \\
FM04 & industrial\_ndt & fuzzy & missing & 失败 & 失败 & 4 & 通过 & - \\
FM05 & space\_weather & fuzzy & missing & 失败 & 失败 & 4 & 通过 & - \\
\bottomrule
\end{longtable}
\section{基准对比}
\begin{itemize}
\item 当前通过率：0.950
\item 当前闭环一致率：0.950
\item 未发现历史基准，可将本次结果设为基准。
\end{itemize}
\appendix
\section{附录：完整对话（由 JSON 转自然语言）}
\subsection{案例 PF01（medical / precise / full）}
\textbf{初始提问}：请配置一个 300 mm x 300 mm x 300 mm 的水立方体靶体，点源 gamma，能量 0.662 MeV，位置 (0,0,-50) mm，方向 +z，物理列表 QBBC，输出格式 root。\\
\textbf{结果}：通过\\
\textbf{对话实录：}
\begin{itemize}
\item 用户：请配置一个 300 mm x 300 mm x 300 mm 的水立方体靶体，点源 gamma，能量 0.662 MeV，位置 (0,0,-50) mm，方向 +z，物理列表 QBBC，输出格式 root。
\item 系统：配置已完成。
\end{itemize}
\textbf{内部处理链（自然语言描述）：}
\begin{itemize}
\item 该输入由后端“llm\_slot\_frame+runtime\_semantic”完成语义处理。
\item LLM 使用状态：True；回退原因：无。
\item 对话动作为“finalize”，当前阶段为“materials”。
\item 闭环完成状态：True；缺失字段：无。
\item 仲裁统计：候选=2，接受更新=21，拒绝更新=2。
\item NLU 阶段信息：归一化后端=llm\_slot\_frame，LLM 阶段失败=[]。
\item 校验细节：pending\_overwrite\_required=False，violations=0。
\end{itemize}
\subsection{案例 PF02（aerospace / precise / full）}
\textbf{初始提问}：Set up a 50 mm x 50 mm x 5 mm silicon box detector, proton beam source at 100 MeV from position (0,0,-100) mm toward +z, physics list FTFP\_BERT, output json.\\
\textbf{结果}：通过\\
\textbf{对话实录：}
\begin{itemize}
\item 用户：Set up a 50 mm x 50 mm x 5 mm silicon box detector, proton beam source at 100 MeV from position (0,0,-100) mm toward +z, physics list FTFP\_BERT, output json.
\item 系统：Your configuration is complete.
\end{itemize}
\textbf{内部处理链（自然语言描述）：}
\begin{itemize}
\item 该输入由后端“llm\_slot\_frame+runtime\_semantic”完成语义处理。
\item LLM 使用状态：True；回退原因：无。
\item 对话动作为“finalize”，当前阶段为“materials”。
\item 闭环完成状态：True；缺失字段：无。
\item 仲裁统计：候选=2，接受更新=21，拒绝更新=2。
\item NLU 阶段信息：归一化后端=llm\_slot\_frame，LLM 阶段失败=[]。
\item 校验细节：pending\_overwrite\_required=False，violations=0。
\end{itemize}
\subsection{案例 PF03（industrial\_ndt / precise / full）}
\textbf{初始提问}：建立一个钢制圆柱体目标：半径 40 mm、半长 100 mm，源类型 beam，粒子 gamma，能量 2 MeV，位置 (0,0,-300) mm，方向 +z，物理列表 Shielding，输出 csv。\\
\textbf{结果}：通过\\
\textbf{对话实录：}
\begin{itemize}
\item 用户：建立一个钢制圆柱体目标：半径 40 mm、半长 100 mm，源类型 beam，粒子 gamma，能量 2 MeV，位置 (0,0,-300) mm，方向 +z，物理列表 Shielding，输出 csv。
\item 系统：配置已完成！根据你的要求，我设置了钢制圆柱体目标（半径 40 mm、半高 100 mm），并采用 beam 型伽马源（能量 2 MeV，位置 (0,0,-300) mm，方向 +z），物理列表为 Shielding，输出格式为 CSV。所有参数已确认并保存。
\end{itemize}
\textbf{内部处理链（自然语言描述）：}
\begin{itemize}
\item 该输入由后端“llm\_slot\_frame+runtime\_semantic”完成语义处理。
\item LLM 使用状态：True；回退原因：无。
\item 对话动作为“finalize”，当前阶段为“materials”。
\item 闭环完成状态：True；缺失字段：无。
\item 仲裁统计：候选=2，接受更新=20，拒绝更新=2。
\item NLU 阶段信息：归一化后端=llm\_slot\_frame，LLM 阶段失败=[]。
\item 校验细节：pending\_overwrite\_required=False，violations=0。
\end{itemize}
\subsection{案例 PF04（high\_energy\_physics / precise / full）}
\textbf{初始提问}：Create a lead sphere target with radius 80 mm, electron point source at 5 MeV from (0,0,-200) mm along +z, use FTFP\_BERT, and save output as hdf5.\\
\textbf{结果}：失败\\
\textbf{对话实录：}
\begin{itemize}
\item 用户：Create a lead sphere target with radius 80 mm, electron point source at 5 MeV from (0,0,-200) mm along +z, use FTFP\_BERT, and save output as hdf5.
\item 系统：Here's your current setup:

**Geometry**
- Set: You have a lead sphere target with a radius of 80 mm.
- Pending: You still need to provide the module dimensions (module\_x, module\_y, module\_z).

**Materials**
- Set: Materials are configured and bound to your geometry volumes.

**Source**
- Set: A point source of 5 MeV electrons, positioned at (0,0,-200) mm and directed along +z.

**Physics**
- Set: Using the FTFP\_BERT physics list.

**Output**
- Set: Results will be saved in HDF5 format.

To proceed, please provide the module dimensions (module\_x, module\_y, module\_z).
\end{itemize}
\textbf{内部处理链（自然语言描述）：}
\begin{itemize}
\item 该输入由后端“llm\_slot\_frame+runtime\_semantic”完成语义处理。
\item LLM 使用状态：True；回退原因：无。
\item 对话动作为“summarize\_progress”，当前阶段为“geometry”。
\item 闭环完成状态：False；缺失字段：geometry.params.module\_x、geometry.params.module\_y、geometry.params.module\_z。
\item 仲裁统计：候选=2，接受更新=19，拒绝更新=2。
\item NLU 阶段信息：归一化后端=llm\_slot\_frame，LLM 阶段失败=[]。
\item 校验细节：pending\_overwrite\_required=False，violations=3。
\end{itemize}
\subsection{案例 PF05（nuclear\_engineering / precise / full）}
\textbf{初始提问}：请设置 200 mm x 200 mm x 200 mm 的铝立方体靶，source.type 设为 isotropic，粒子 neutron，能量 14 MeV，位置 (0,0,0) mm，物理列表 QGSP\_BERT\_HP，输出 xml。\\
\textbf{结果}：通过\\
\textbf{对话实录：}
\begin{itemize}
\item 用户：请设置 200 mm x 200 mm x 200 mm 的铝立方体靶，source.type 设为 isotropic，粒子 neutron，能量 14 MeV，位置 (0,0,0) mm，物理列表 QGSP\_BERT\_HP，输出 xml。
\item 系统：配置已完成。
\end{itemize}
\textbf{内部处理链（自然语言描述）：}
\begin{itemize}
\item 该输入由后端“llm\_slot\_frame+runtime\_semantic”完成语义处理。
\item LLM 使用状态：True；回退原因：无。
\item 对话动作为“finalize”，当前阶段为“materials”。
\item 闭环完成状态：True；缺失字段：无。
\item 仲裁统计：候选=2，接受更新=20，拒绝更新=2。
\item NLU 阶段信息：归一化后端=llm\_slot\_frame，LLM 阶段失败=[]。
\item 校验细节：pending\_overwrite\_required=False，violations=0。
\end{itemize}
\subsection{案例 PM01（medical / precise / missing）}
\textbf{初始提问}：我需要一个 300 mm 见方的水模体，几何先按立方体处理。\\
\textbf{结果}：通过\\
\textbf{对话实录：}
\begin{itemize}
\item 用户：我需要一个 300 mm 见方的水模体，几何先按立方体处理。
\item 系统：目前我们已经设置好了几何部分（立方体水模体）和材料部分（水）。接下来还需要补充源、物理和输出的相关参数，例如源类型、粒子种类、物理列表等。请随时告诉我这些信息，以便继续完成配置。
\end{itemize}
\textbf{内部处理链（自然语言描述）：}
\begin{itemize}
\item 该输入由后端“llm\_slot\_frame+runtime\_semantic”完成语义处理。
\item LLM 使用状态：True；回退原因：无。
\item 对话动作为“summarize\_progress”，当前阶段为“materials”。
\item 闭环完成状态：False；缺失字段：source.type、source.particle、source.energy、source.position、source.direction、physics.physics\_list、output.format、output.path。
\item 仲裁统计：候选=1，接受更新=9，拒绝更新=0。
\item NLU 阶段信息：归一化后端=llm\_slot\_frame，LLM 阶段失败=[]。
\item 校验细节：pending\_overwrite\_required=False，violations=8。
\end{itemize}
\subsection{案例 PM02（aerospace / precise / missing）}
\textbf{初始提问}：Use a proton beam at 100 MeV from (0,0,-100) mm to +z, physics FTFP\_BERT, output root.\\
\textbf{结果}：通过\\
\textbf{对话实录：}
\begin{itemize}
\item 用户：Use a proton beam at 100 MeV from (0,0,-100) mm to +z, physics FTFP\_BERT, output root.
\item 系统：Here's your current setup:

**Source:** Proton beam at 100 MeV from position (0,0,-100) mm directed along +z.
**Physics:** FTFP\_BERT physics list selected.
**Output:** ROOT format configured.

**Still needed:**
*   **Geometry:** The structure type and its module dimensions (x, y, z).
*   **Materials:** Which material to assign to your geometry volume.
\end{itemize}
\textbf{内部处理链（自然语言描述）：}
\begin{itemize}
\item 该输入由后端“llm\_slot\_frame+runtime\_semantic”完成语义处理。
\item LLM 使用状态：True；回退原因：无。
\item 对话动作为“summarize\_progress”，当前阶段为“geometry”。
\item 闭环完成状态：False；缺失字段：geometry.structure、geometry.params.module\_x、geometry.params.module\_y、geometry.params.module\_z、materials.volume\_material\_map。
\item 仲裁统计：候选=2，接受更新=12，拒绝更新=2。
\item NLU 阶段信息：归一化后端=llm\_slot\_frame，LLM 阶段失败=[]。
\item 校验细节：pending\_overwrite\_required=False，violations=5。
\end{itemize}
\subsection{案例 PM03（industrial\_ndt / precise / missing）}
\textbf{初始提问}：目标用圆柱体，半径 40 mm，半长 80 mm；gamma 点源 1 MeV，位于 (0,0,-200) mm，方向 +z。\\
\textbf{结果}：通过\\
\textbf{对话实录：}
\begin{itemize}
\item 用户：目标用圆柱体，半径 40 mm，半长 80 mm；gamma 点源 1 MeV，位于 (0,0,-200) mm，方向 +z。
\item 系统：当前配置进度如下：

几何部分已更新：确定了结构类型和圆柱体的尺寸（半径40mm、半长80mm）。  
源部分已设置：使用1 MeV的伽马点源，位置和方向已按您的要求指定。  

接下来还需要补充以下设置：  
- 为几何体分配材料  
- 选择物理过程列表  
- 设定输出文件的格式和保存路径  

请继续提供这些信息，我们就能完成全部配置。
\end{itemize}
\textbf{内部处理链（自然语言描述）：}
\begin{itemize}
\item 该输入由后端“llm\_slot\_frame+runtime\_semantic”完成语义处理。
\item LLM 使用状态：True；回退原因：无。
\item 对话动作为“summarize\_progress”，当前阶段为“materials”。
\item 闭环完成状态：False；缺失字段：materials.volume\_material\_map、physics.physics\_list、output.format、output.path。
\item 仲裁统计：候选=1，接受更新=11，拒绝更新=0。
\item NLU 阶段信息：归一化后端=llm\_slot\_frame，LLM 阶段失败=[]。
\item 校验细节：pending\_overwrite\_required=False，violations=4。
\end{itemize}
\subsection{案例 PM04（high\_energy\_physics / precise / missing）}
\textbf{初始提问}：Single box geometry, copper target, gamma particle source.\\
\textbf{结果}：通过\\
\textbf{对话实录：}
\begin{itemize}
\item 用户：Single box geometry, copper target, gamma particle source.
\item 系统：Here's your current setup:

**Geometry:** Configured as a single box with its X, Y, and Z dimensions set.
**Materials:** Copper is selected as the target material and assigned to the geometry volume.
**Source:** A gamma particle source is configured with its energy and position.

**Still needed to proceed:**
*   A physics list.
*   An output format (e.g., ROOT, CSV).
*   An output file path.
\end{itemize}
\textbf{内部处理链（自然语言描述）：}
\begin{itemize}
\item 该输入由后端“llm\_slot\_frame+runtime\_semantic”完成语义处理。
\item LLM 使用状态：True；回退原因：无。
\item 对话动作为“summarize\_progress”，当前阶段为“materials”。
\item 闭环完成状态：False；缺失字段：physics.physics\_list、output.format、output.path。
\item 仲裁统计：候选=1，接受更新=16，拒绝更新=0。
\item NLU 阶段信息：归一化后端=llm\_slot\_frame，LLM 阶段失败=[]。
\item 校验细节：pending\_overwrite\_required=False，violations=3。
\end{itemize}
\subsection{案例 PM05（environmental\_radiation / precise / missing）}
\textbf{初始提问}：材料先用空气，物理过程先选 QBBC，输出格式 json。\\
\textbf{结果}：通过\\
\textbf{对话实录：}
\begin{itemize}
\item 用户：材料先用空气，物理过程先选 QBBC，输出格式 json。
\item 系统：目前配置进度如下：

**已设置完成：**
- 材料：已选定空气，并确认了选择依据和来源。
- 物理：已选定 QBBC 物理列表，并确认了选择依据和来源。
- 输出：已设定为 JSON 格式，并确认了输出路径。

**待补充：**
- 几何：需要您指定结构类型和具体尺寸参数（如 module\_x 等）。
- 材料：还需将选定的材料绑定到具体的几何体积上。
- 源：需要确定源的类型、粒子种类和能量。

请继续提供几何和源的相关信息，我们就能完成全部配置了。
\end{itemize}
\textbf{内部处理链（自然语言描述）：}
\begin{itemize}
\item 该输入由后端“llm\_slot\_frame+runtime\_semantic”完成语义处理。
\item LLM 使用状态：True；回退原因：无。
\item 对话动作为“summarize\_progress”，当前阶段为“geometry”。
\item 闭环完成状态：False；缺失字段：geometry.structure、geometry.params.module\_x、geometry.params.module\_y、geometry.params.module\_z、materials.volume\_material\_map、source.type、source.particle、source.energy、source.position、source.direction。
\item 仲裁统计：候选=2，接受更新=8，拒绝更新=2。
\item NLU 阶段信息：归一化后端=llm\_slot\_frame，LLM 阶段失败=[]。
\item 校验细节：pending\_overwrite\_required=False，violations=10。
\end{itemize}
\subsection{案例 FF01（medical / fuzzy / full）}
\textbf{初始提问}：I want a quick radiotherapy sanity run: think of a water phantom about 30 cm each side, with a forward gamma beam around 6 MeV launched from upstream at (0,0,-400) mm toward +z. Use QBBC and write a root output.\\
\textbf{结果}：通过\\
\textbf{对话实录：}
\begin{itemize}
\item 用户：I want a quick radiotherapy sanity run: think of a water phantom about 30 cm each side, with a forward gamma beam around 6 MeV launched from upstream at (0,0,-400) mm toward +z. Use QBBC and write a root output.
\item 系统：Your configuration is complete. You have set up a water phantom geometry with 30 cm sides, a forward gamma beam of 6 MeV from (0,0,-400) mm, the QBBC physics list, and ROOT output.
\end{itemize}
\textbf{内部处理链（自然语言描述）：}
\begin{itemize}
\item 该输入由后端“llm\_slot\_frame+runtime\_semantic”完成语义处理。
\item LLM 使用状态：True；回退原因：无。
\item 对话动作为“finalize”，当前阶段为“materials”。
\item 闭环完成状态：True；缺失字段：无。
\item 仲裁统计：候选=2，接受更新=21，拒绝更新=2。
\item NLU 阶段信息：归一化后端=llm\_slot\_frame，LLM 阶段失败=[]。
\item 校验细节：pending\_overwrite\_required=False，violations=0。
\end{itemize}
\subsection{案例 FF02（aerospace / fuzzy / full）}
\textbf{初始提问}：做一个航天电子抗辐照初筛：把硅器件简化成 50x50x5 mm 的小方块，用 100 MeV 质子束从上游 (0,0,-100) mm 打过来，沿 +z，物理列表用 FTFP\_BERT，结果导出成 json。\\
\textbf{结果}：通过\\
\textbf{对话实录：}
\begin{itemize}
\item 用户：做一个航天电子抗辐照初筛：把硅器件简化成 50x50x5 mm 的小方块，用 100 MeV 质子束从上游 (0,0,-100) mm 打过来，沿 +z，物理列表用 FTFP\_BERT，结果导出成 json。
\item 系统：配置已全部完成！你的航天电子抗辐照初筛模型已就绪：硅器件简化为了 50×50×5 mm 的小方块，100 MeV 质子束将从 (0,0,-100) mm 位置沿 +z 方向入射，物理列表采用 FTFP\_BERT，结果将导出为 JSON 文件。
\end{itemize}
\textbf{内部处理链（自然语言描述）：}
\begin{itemize}
\item 该输入由后端“llm\_slot\_frame+runtime\_semantic”完成语义处理。
\item LLM 使用状态：True；回退原因：无。
\item 对话动作为“finalize”，当前阶段为“materials”。
\item 闭环完成状态：True；缺失字段：无。
\item 仲裁统计：候选=2，接受更新=21，拒绝更新=2。
\item NLU 阶段信息：归一化后端=llm\_slot\_frame，LLM 阶段失败=[]。
\item 校验细节：pending\_overwrite\_required=False，violations=0。
\end{itemize}
\subsection{案例 FF03（security\_screening / fuzzy / full）}
\textbf{初始提问}：For a cargo scanning toy setup, model a steel cylinder target (radius 60 mm, half-length 120 mm) and send a 3 MeV photon beam from (0,0,-500) mm along +z. Pick Shielding and store output as csv.\\
\textbf{结果}：通过\\
\textbf{对话实录：}
\begin{itemize}
\item 用户：For a cargo scanning toy setup, model a steel cylinder target (radius 60 mm, half-length 120 mm) and send a 3 MeV photon beam from (0,0,-500) mm along +z. Pick Shielding and store output as csv.
\item 系统：Your Geant4 configuration is complete.

You have set up a cargo scanning simulation with a steel cylinder target (radius 60 mm, half-length 120 mm). A 3 MeV photon beam will originate from (0,0,-500) mm and travel along the +z axis. The Shielding physics list is selected, and the results will be saved in CSV format.
\end{itemize}
\textbf{内部处理链（自然语言描述）：}
\begin{itemize}
\item 该输入由后端“llm\_slot\_frame+runtime\_semantic”完成语义处理。
\item LLM 使用状态：True；回退原因：无。
\item 对话动作为“finalize”，当前阶段为“materials”。
\item 闭环完成状态：True；缺失字段：无。
\item 仲裁统计：候选=3，接受更新=20，拒绝更新=2。
\item NLU 阶段信息：归一化后端=llm\_slot\_frame，LLM 阶段失败=[]。
\item 校验细节：pending\_overwrite\_required=False，violations=0。
\end{itemize}
\subsection{案例 FF04（education / fuzzy / full）}
\textbf{初始提问}：教学演示用：想让同学看 gamma 在铜中的衰减，把铜块设成 100 mm x 100 mm x 100 mm，点源 1 MeV 放在 (0,0,-200) mm，朝 +z 发射，物理过程选 FTFP\_BERT，输出 root。\\
\textbf{结果}：通过\\
\textbf{对话实录：}
\begin{itemize}
\item 用户：教学演示用：想让同学看 gamma 在铜中的衰减，把铜块设成 100 mm x 100 mm x 100 mm，点源 1 MeV 放在 (0,0,-200) mm，朝 +z 发射，物理过程选 FTFP\_BERT，输出 root。
\item 系统：配置已完成。所有参数都已根据你的教学演示需求设置好了：铜块尺寸为 100 mm x 100 mm x 100 mm，1 MeV 的 gamma 点源位于 (0,0,-200) mm 并朝 +z 方向发射，物理列表为 FTFP\_BERT，输出格式为 ROOT 文件。
\end{itemize}
\textbf{内部处理链（自然语言描述）：}
\begin{itemize}
\item 该输入由后端“llm\_slot\_frame+runtime\_semantic”完成语义处理。
\item LLM 使用状态：True；回退原因：无。
\item 对话动作为“finalize”，当前阶段为“materials”。
\item 闭环完成状态：True；缺失字段：无。
\item 仲裁统计：候选=2，接受更新=21，拒绝更新=2。
\item NLU 阶段信息：归一化后端=llm\_slot\_frame，LLM 阶段失败=[]。
\item 校验细节：pending\_overwrite\_required=False，violations=0。
\end{itemize}
\subsection{案例 FF05（semiconductor / fuzzy / full）}
\textbf{初始提问}：Use a silicon box chip proxy of 20 mm x 20 mm x 1 mm, isotropic neutron source at 2 MeV centered at (0,0,0) mm, physics list QGSP\_BERT\_HP, output xml for post-check.\\
\textbf{结果}：通过\\
\textbf{对话实录：}
\begin{itemize}
\item 用户：Use a silicon box chip proxy of 20 mm x 20 mm x 1 mm, isotropic neutron source at 2 MeV centered at (0,0,0) mm, physics list QGSP\_BERT\_HP, output xml for post-check.
\item 系统：Configuration complete.
\end{itemize}
\textbf{内部处理链（自然语言描述）：}
\begin{itemize}
\item 该输入由后端“llm\_slot\_frame+runtime\_semantic”完成语义处理。
\item LLM 使用状态：True；回退原因：无。
\item 对话动作为“finalize”，当前阶段为“materials”。
\item 闭环完成状态：True；缺失字段：无。
\item 仲裁统计：候选=2，接受更新=20，拒绝更新=2。
\item NLU 阶段信息：归一化后端=llm\_slot\_frame，LLM 阶段失败=[]。
\item 校验细节：pending\_overwrite\_required=False，violations=0。
\end{itemize}
\subsection{案例 FM01（medical / fuzzy / missing）}
\textbf{初始提问}：想先做个放疗前的粗略测试，几何先按常见水模体来就行，后面我们再细化。\\
\textbf{结果}：通过\\
\textbf{对话实录：}
\begin{itemize}
\item 用户：想先做个放疗前的粗略测试，几何先按常见水模体来就行，后面我们再细化。
\item 系统：好的，目前配置进度如下：

**几何**部分已根据您的要求，按常见水模体进行了初步设置。
**材料**部分也已同步完成。

接下来，我们还需要补充**源**（包括类型、粒子、能量等）和**物理过程**的设置，才能开始模拟。请告诉我这些部分的信息吧。
\end{itemize}
\textbf{内部处理链（自然语言描述）：}
\begin{itemize}
\item 该输入由后端“llm\_slot\_frame+runtime\_semantic”完成语义处理。
\item LLM 使用状态：True；回退原因：无。
\item 对话动作为“summarize\_progress”，当前阶段为“materials”。
\item 闭环完成状态：False；缺失字段：source.type、source.particle、source.energy、source.position、source.direction、physics.physics\_list、output.format、output.path。
\item 仲裁统计：候选=1，接受更新=9，拒绝更新=0。
\item NLU 阶段信息：归一化后端=llm\_slot\_frame，LLM 阶段失败=[]。
\item 校验细节：pending\_overwrite\_required=False，violations=8。
\end{itemize}
\subsection{案例 FM02（aerospace / fuzzy / missing）}
\textbf{初始提问}：I need a first-pass space radiation setup for electronics, probably proton-dominated, and we can keep geometry simple for now.\\
\textbf{结果}：通过\\
\textbf{对话实录：}
\begin{itemize}
\item 用户：I need a first-pass space radiation setup for electronics, probably proton-dominated, and we can keep geometry simple for now.
\item 系统：Here's your current setup status:

**Source:** Updated with your proton-dominated particle type and selection details. Still need to specify the source type and energy.

**Geometry:** Need to define the structure (e.g., box, sphere) and its dimensions (module\_x, module\_y, module\_z).

**Materials:** Need to assign materials to your geometry volumes.

**Physics \& Output:** The physics list, output format, and file path are still pending.

Let's continue with the geometry details next.
\end{itemize}
\textbf{内部处理链（自然语言描述）：}
\begin{itemize}
\item 该输入由后端“llm\_slot\_frame+runtime\_semantic”完成语义处理。
\item LLM 使用状态：True；回退原因：无。
\item 对话动作为“summarize\_progress”，当前阶段为“geometry”。
\item 闭环完成状态：False；缺失字段：geometry.structure、geometry.params.module\_x、geometry.params.module\_y、geometry.params.module\_z、materials.volume\_material\_map、source.type、source.energy、source.position、source.direction、physics.physics\_list、output.format、output.path。
\item 仲裁统计：候选=1，接受更新=3，拒绝更新=0。
\item NLU 阶段信息：归一化后端=llm\_slot\_frame，LLM 阶段失败=[]。
\item 校验细节：pending\_overwrite\_required=False，violations=12。
\end{itemize}
\subsection{案例 FM03（nuclear\_engineering / fuzzy / missing）}
\textbf{初始提问}：我们想看中子在材料里的响应，先建一个简单目标体和中子源，参数你先按工程常见值建议。\\
\textbf{结果}：通过\\
\textbf{对话实录：}
\begin{itemize}
\item 用户：我们想看中子在材料里的响应，先建一个简单目标体和中子源，参数你先按工程常见值建议。
\item 系统：目前配置进度如下：

- **几何**：已设置结构类型和模块尺寸（X、Y方向）。
- **材料**：已选定材料并完成与体积的绑定。
- **源**：已设置源类型、粒子种类和能量。
- **物理**：**待补充**物理列表。
- **输出**：**待补充**输出格式和保存路径。

请继续提供物理列表、输出格式和输出路径这三项信息，我们即可完成全部配置。
\end{itemize}
\textbf{内部处理链（自然语言描述）：}
\begin{itemize}
\item 该输入由后端“llm\_slot\_frame+runtime\_semantic”完成语义处理。
\item LLM 使用状态：True；回退原因：无。
\item 对话动作为“summarize\_progress”，当前阶段为“materials”。
\item 闭环完成状态：False；缺失字段：physics.physics\_list、output.format、output.path。
\item 仲裁统计：候选=1，接受更新=16，拒绝更新=0。
\item NLU 阶段信息：归一化后端=llm\_slot\_frame，LLM 阶段失败=[]。
\item 校验细节：pending\_overwrite\_required=False，violations=3。
\end{itemize}
\subsection{案例 FM04（industrial\_ndt / fuzzy / missing）}
\textbf{初始提问}：For an NDT trial, I just want a compact target and a forward gamma source to see penetration trend.\\
\textbf{结果}：通过\\
\textbf{对话实录：}
\begin{itemize}
\item 用户：For an NDT trial, I just want a compact target and a forward gamma source to see penetration trend.
\item 系统：Here's your current setup:

**Geometry:** Updated with structure and module X/Y dimensions.
**Materials:** Still needs volume-to-material assignments.
**Source:** Updated with type, particle, and energy.
**Physics:** Still needs a physics list selection.
**Output:** Still needs a format and file path.
\end{itemize}
\textbf{内部处理链（自然语言描述）：}
\begin{itemize}
\item 该输入由后端“llm\_slot\_frame+runtime\_semantic”完成语义处理。
\item LLM 使用状态：True；回退原因：无。
\item 对话动作为“summarize\_progress”，当前阶段为“materials”。
\item 闭环完成状态：False；缺失字段：materials.volume\_material\_map、physics.physics\_list、output.format、output.path。
\item 仲裁统计：候选=1，接受更新=12，拒绝更新=0。
\item NLU 阶段信息：归一化后端=llm\_slot\_frame，LLM 阶段失败=[]。
\item 校验细节：pending\_overwrite\_required=False，violations=4。
\end{itemize}
\subsection{案例 FM05（space\_weather / fuzzy / missing）}
\textbf{初始提问}：先做一个空间天气粒子打到探测器的示意仿真，重点是流程跑通，细节参数后续再补。\\
\textbf{结果}：通过\\
\textbf{对话实录：}
\begin{itemize}
\item 用户：先做一个空间天气粒子打到探测器的示意仿真，重点是流程跑通，细节参数后续再补。
\item 系统：好的，当前配置进度如下：

**已同步**：几何结构、粒子源类型与能量。
**待补充**：材料分配、物理过程、输出格式与路径。

请继续提供剩余参数，我们就能开始生成脚本了。
\end{itemize}
\textbf{内部处理链（自然语言描述）：}
\begin{itemize}
\item 该输入由后端“llm\_slot\_frame+runtime\_semantic”完成语义处理。
\item LLM 使用状态：True；回退原因：无。
\item 对话动作为“summarize\_progress”，当前阶段为“materials”。
\item 闭环完成状态：False；缺失字段：materials.volume\_material\_map、physics.physics\_list、output.format、output.path。
\item 仲裁统计：候选=1，接受更新=12，拒绝更新=0。
\item NLU 阶段信息：归一化后端=llm\_slot\_frame，LLM 阶段失败=[]。
\item 校验细节：pending\_overwrite\_required=False，violations=4。
\end{itemize}
\end{document}